% Part: sets-functions-relations
% Chapter: arithmetization
% Section: cauchy

\documentclass[../../../include/open-logic-section]{subfiles}

\begin{document}

\olfileid[ar]{sfr}{arith}{cauchy}

\olsection{ملحق: بناء الحقيقية بمتتاليات كوشي}

كان بناء الأعداد الحقيقية في \olref[cuts]{sec} بقطوع ديديكند. ونبين هنا بناء آخر، أصله تعريف كوشي لما يسمى اليوم متتالية كوشي. أما اتخاذ هذا التعريف سبيلًا إلى \emph{بناء} الحقيقية فمن عمل مؤلفين آخرين في القرن التاسع عشر، منهم فايرشتراس وهاينه وميري وكانتور. وفي تاريخ ذلك عرض حسن عند \citeauthor{OConnorRobertson:RN}، \citeyear{OConnorRobertson:RN}.

وقبل بلوغ أولئك، لننظر في فكرة سيمون ستيفن (1548--1620). فقد أدرك أن كل عدد حقيقي يمكن النظر إليه من جهة تمثيله العشري؛ حتى غير النسبي، مثل $\sqrt{2}$، له تمثيل عشري معلوم أولُه:
\[
	1.41421356237\ldots
\]
وتمثيل هذه الكسور العشرية في نظرية المجموعات يسير: نجعلها دوال $d \colon \Nat \to \Nat$، فتكون $d(n)$ الخانة العشرية ذات الرتبة $n$ التي ننظر فيها، ونقيد قيم الخانات بالأرقام العشرية. ثم نحتاج إلى تعديل صغير لحفظ الجزء السابق للفاصلة العشرية، وهو هنا $1$، وإلى تعديل آخر بعلاقة تكافؤ يضمن، مثلًا، $0.999\ldots = 1$. ولا يصعب بعد ذلك إقامة بناء صارم تمامًا للأعداد الحقيقية على طريقة ستيفن في نظرية المجموعات.

وليس تفصيل أمر ستيفن غرضنا (وللزيادة انظر \citealt{KatzKatz2012})، لكن فكرته تقود إلى معنى قريب. فعوضًا عن النظر مباشرة في التمثيل العشري لـ$\sqrt{2}$، نأخذ \emph{متتالية} من التقريبات النسبية لـ$\sqrt{2}$ تزداد دقة شيئًا فشيئًا:
\[
1, 1.4, 1.414, 1.4142, 1.41421,\ldots
\]
ومن هذا النظر إلى الحقيقية بواسطة «تقريبات تتحسن تدريجيًّا» تنشأ سلسلة ملاحظات، نظير ما استعملناه لبناء $\Rat$ من $\Int$، أو $\Int$ من $\Nat$:
\begin{enumerate}
\item كل عدد حقيقي له تمثيل عشري، وقد يكون غير منته.
\item ما يحمله التمثيل العشري من معلومات، وإن لم ينته، يمكن أن تحمله متتالية من الأعداد النسبية.
\item ومتتالية الأعداد النسبية دالة من $\Nat$ إلى $\Rat$؛ نجعل $f(n)$ العدد النسبي ذا الرتبة $n$ فيها.
\end{enumerate}
على أن \emph{مطلق} الدالة من $\Nat$ إلى $\Rat$ لا يحدد عددًا حقيقيًّا. وشاهد ذلك الدالة
\[
	f(n) =\begin{cases}
			1 & \text{إذا كان }n\text{ فرديًا}\\
			0 &\text{إذا كان }n\text{ زوجيًا}
		\end{cases}
\]
فالمتتالية $0,1,0,1,0,1,0,\ldots$ لا تتجه إلى عدد حقيقي بعينه. فلا بد أن نقصر النظر على ما يقترب من عدد حقيقي، أي على متتاليات ذات \emph{نهاية}.

وقد مر معنى النهاية في \olref[his][set][limits]{sec}. لكن لا يسوغ أن نأخذ التعريف السابق \emph{بعينه}؛ إذ كان في عبارة «$(\forall \epsilon>0)$» تكميم ضمني على الأعداد الحقيقية، وكانت الدوال المنظور في نهاياتها على الأعداد الحقيقية. فاستعماله هنا يفترض توافر ما نريد \emph{بناءه}. غير أن معنى قريبًا من النهاية يغنينا عنه.
\begin{defn}\ollabel{def:CauchySequence}
الدالة $f: \Nat \to \Rat$ \emph{متتالية كوشي} إذا وفقط إذا تحقق، لكل $\epsilon \in \Rat$ موجب، أن $(\exists \ell \in
\Nat)(\forall m, n > \ell)|f(m) - f(n)| < \epsilon$.
\end{defn}

فالمعنى العام قريب مما تقدم: مهما طلبنا قدرًا من الدقة يقاس بـ~$\epsilon$، وجدنا موضعًا إذا تجاوزناه---أي أخذنا المدخل أكبر من~$\ell$---حصل المقصود. ويتبين من هذا أن متتاليتنا $1$، و$1.4$، و$1.414$، و$1.4142$، و$1.41421$\ldots تتجه إلى نهاية؛ فلتقريب $\sqrt{2}$ بخطأ لا يجاوز $\nicefrac{1}{10^{n}}$ يكفي أن نأخذ أي حد بعد الحد ذي الرتبة $n$.

فيتبادر أن العدد الحقيقي \emph{هو} متتالية كوشي كيف اتفقت. غير أن هذا، كما في بناء $\Int$ و$\Rat$، لا يفي بالمقصود؛ فقد تدل متتاليات مختلفة على العدد الحقيقي الواحد. ومثال ذلك أن نأخذ متتالية كوشي~$f$، ونعين $g$ مثل~$f$ في جميع القيم سوى $g(0)\neq f(0)$. وهما تتفقان فيما بعد الحد الأول، فنريد آخر الأمر أن نجعل لهما النهاية نفسها، على المعنى المقصود في \olref{def:CauchySequence}؛ ومن ثم تعدان طريقتين «لتعريف» عدد حقيقي واحد. فلا بد من معاملتهما من هذه الجهة معاملة العدد نفسه.

فنحتاج مرة أخرى إلى علاقة تكافؤ على متتاليات كوشي، ثم نطابق الأعداد الحقيقية مع فئاتها، لا مع العلاقات نفسها. ونقدم لذلك معنى الدالة الآيلة إلى $0$. فإذا كانت $h : \Nat \to \Rat$، قلنا \emph{إن $h$ تؤول إلى $0$} إذا وفقط إذا تحقق لكل $\epsilon \in
\Rat$ موجب أن $(\exists \ell \in \Nat)(\forall n > \ell)|h(n)| <
\epsilon$.\footnote{قابل هذا بتعريف $\lim_{x
\mathord{\rightarrow}\infty}f(x) = 0$ في \olref[his][set][limits]{sec}.} وإذا كانت $f$ و$g$ دالتين من النوع $\Nat \to \Rat$، وضعنا $(f-g)(n) = f(n) - g(n)$. ثم نعرف
\[
	f \Realequiv g \text{ إذا وفقط إذا كانت الدالة $(f-g)$ تؤول إلى $0$}.
\]
ويلزم إثبات أن $\Realequiv$ علاقة تكافؤ، وهو صحيح. فبعده يصح أن نعرف الأعداد الحقيقية بأنها فئات التكافؤ تحت $\Realequiv$ لمتتاليات كوشي كلها من النوع $\Nat \to \Rat$.

\begin{prob}
ضع $f(n) = 0$ لكل $n$، وضع $g(n) = \frac{1}{(n+1)^2}$. أثبت أن الدالتين متتاليتا كوشي، بل إنهما تؤولان إلى $0$، فيتحقق أيضًا $f \sim_\Real g$.
\end{prob}

ونكتب، على عادتنا، $\equivrep{f}{\Realequiv}$ لفئة التكافؤ التي تضم $f$ بوصفها !!a{element} منها. ونحذف الرمز السفلي فيما يلي تخفيفًا للقراءة، فنكتب $\equivrep{f}{}$. ونضع لكل $q \in \Rat$ العدد $q_{\Real} = \equivrep{c_{q}}{}$، حيث $c_{q}$ الدالة الثابتة التي تحقق $c_q(n) = q$ لكل $n \in \Nat$. ثم نعين العلاقات والعمليات الأساسية، ومن ذلك
\begin{align*}
	\equivrep{f}{} + \equivrep{g}{} &= 	\equivrep{(f + g)}{} \\
	\equivrep{f}{} \times 	\equivrep{g}{} &= \equivrep{(f \times g)}{} 
\end{align*}
ونريد هنا $(f + g)(n) = f(n) + g(n)$ و$(f \times g)(n) = f(n) \times
g(n)$. ويلزم أن نثبت أن $(f + g)$ و$(f-g)$ و$(f\times g)$ كلها متتاليات كوشي متى كانت $f$ و$g$ كذلك؛ وهذه أحكام صحيحة نترك إثباتها للقارئ.

ونعين أخيرًا الترتيب: تكون $\equivrep{f}{}$ \emph{موجبة} إذا وفقط إذا اجتمع الشرطان $\equivrep{f}{}\neq 0_\Real$ و$(\exists \ell \in \Nat)(\forall n > \ell)0 < f(n)$. ثم نضع $\equivrep{f}{} < \equivrep{g}{}$ إذا وفقط إذا كانت $\equivrep{(g - f)}{}$ موجبة. ولا بد من إثبات أن هذا التعيين لا يتغير بتغير الدالة «الممثلة» لفئة التكافؤ.

\textbf{تنبيه تصحيحي على الأصل.} قورنت فئة التكافؤ هنا بالصفر الحقيقي
$0_\Real$، لا بالصفر النسبي $0_\Rat$؛ فذلك يصون نوع طرفي المقارنة.
%\begin{align*}
%	\equivrep{f}{} \leq \equivrep{g}{} \text{ iff }&\equivrep{f}{} = \equivrep{g}{} \text{ or }\\
%	&(\exists \ell \in \Nat)(\forall n \in \Nat)(\ell < n \lif f(n) \leq g(n))
%\end{align*}
فإذا تم ذلك، سهل إثبات أن التعريفات تحقق الخواص الجبرية المطلوبة:
\begin{thm}\ollabel{thm:cauchyorderedfield}
فئات تكافؤ متتاليات كوشي حقل مرتب.
\end{thm}

\begin{proof}
يترك البرهان تمرينًا.
\end{proof}

\begin{prob}
أثبت أن فئات تكافؤ متتاليات كوشي حقل مرتب.
\end{prob}

وأما اكتمال الأعداد الحقيقية بهذا البناء فإثباته أشق، ولذلك نسوق برهانه.

\begin{thm}
كل مجموعة غير خالية من متتاليات كوشي، مأخوذة بفئات تكافئها، إذا كان لها حد علوي فلها أصغر حد علوي.
\end{thm}

\begin{proof}[مخطط البرهان]
خذ $S$ مجموعة غير خالية من متتاليات كوشي لها حد علوي، مع فهم المتتاليات من جهة فئات تكافئها. فيوجد $p \in \Rat$ يكون $p_{\Real}$ حدًّا علويًّا لـ$S$. وخذ $r \in S$؛ فيوجد $q \in \Rat$ يحقق $q_{\Real} < \equivrep{r}{}$. فإذا وجد لـ$S$ أصغر حد علوي، وقع بين $q_\Real$ و$p_\Real$، مع إدخال الطرفين.

\textbf{تنبيه تصحيحي على الأصل.} أُظهر قوسا الفئة في المقارنة
$q_\Real < \equivrep{r}{}$؛ لأن $r$ ممثل من متتاليات كوشي، والترتيب
معرّف على فئات التكافؤ.

فنضيق المسافة إلى أصغر حد علوي من الجهتين معًا. ونعين لذلك دالتين $f, g \colon
\Nat \to \Rat$، لتقترب $f$ منه من أعلى،\ وتقترب $g$ منه من أسفل. ونبدأ بوضع
\begin{align*}
	f(0) &= p \\
	g(0) &= q
\end{align*}
ثم نضع $a_n = \frac{f(n) + g(n)}{2}$، ونعرف:\footnote{هذا تعريف عودي؛ ولم نقدم \emph{إلى الآن} ما يسوغ قبول التعريفات العودية.}
\begin{align*}
	f(n+1) &=
	\begin{cases}
		a_n &\text{إذا كان }(\forall h \in S)\equivrep{h}{} \leq (a_n)_\Real\\
		f(n)&\text{فيما عدا ذلك}
	\end{cases}\\
	g(n+1) &=
	\begin{cases}
		a_n &\text{إذا كان }(\exists h \in S)\equivrep{h}{} \geq (a_n)_\Real\\
	 	g(n) &\text{فيما عدا ذلك}
	\end{cases}
\end{align*}
والدالتان $f$ و$g$ متتاليتا كوشي؛ والتحقق من ذلك يسير نسبيًّا، فنتركه تمرينًا. ثم إن $(f-g)$ تؤول إلى $0$؛ لأن المسافة بين $f$ و$g$ لا تزيد في كل خطوة على نصف سابقتها. ومنه $\equivrep{f}{} = \equivrep{g}{}$.

نثبت أولًا أن $\equivrep{f}{}$ حد علوي لـ$S$، أي\ 
$(\forall h \in S)\equivrep{h}{} \leq \equivrep{f}{}$.
ونستعمل في ذلك \olref{thm:cauchyorderedfield}. خذ $h \in S$، وافرض للخلف أن $\equivrep{f}{} < \equivrep{h}{}$؛ فيكون $0_\Real < \equivrep{(h-f)}{}$. وحيث كانت $f$ متتالية كوشي متناقصة رتيبًا، وجد $n \in \Nat$ يحقق $\equivrep{(c_{f(n)} - f)}{} < \equivrep{(h-f)}{}$. فينتج
\[
	(f(n))_\Real = \equivrep{c_{f(n)}}{} < \equivrep{f}{} + \equivrep{(h-f)}{} = \equivrep{h}{},
\]
وهذا يناقض ما يضمنه البناء من $\equivrep{h}{} \leq (f(n))_\Real$.

ثم نثبت أن $\equivrep{f}{} = \equivrep{g}{}$ هو \emph{أصغر} حد علوي لـ$S$. خذ $j$ متتالية كوشي كيف اتفقت، وافرض $\equivrep{j}{} <
\equivrep{g}{}$. فبمثل الاستدلال السابق، مع استعمال كون $g$ \emph{متزايدة}، يوجد $n \in \Nat$ يحقق $\equivrep{j}{} <
(g(n))_\Real$. والبناء يضمن وجود $h \in S$ بحيث $(g(n))_\Real \leq \equivrep{h}{}$. فيلزم $\equivrep{j}{} <
\equivrep{h}{}$، فلا تكون $\equivrep{j}{}$ حدًّا علويًّا لـ$S$، وهو المطلوب.
\end{proof}

\end{document}
